\documentclass[../main.tex]{subfiles}
\graphicspath{{\subfix{../images/}}}
\begin{document}
	\subsection{Carica elettrone}
	La carica di un elettrone ($e^-$) è: $-1.6 \cdot 10^{-19} C$ \\
	L'unita di misura utilizzata è il Coulomb.
	
	\subsection{Prima legge di Ohm}
	\begin{equation}
		\Delta V = R \cdot I \;\; [V]
	\end{equation}
	$\Delta V$ è la \textbf{differenza di potenziale} che si misura in Volt $[V]$ \\
	$R$ è la \textbf{resistenza} che si misura in Ohm $[\Omega]$ \\
	$I$ è la \textbf{corrente elettrica} che si misura in Ampere $[A]$\\\\
	La differenza di potenziale \textbf{può anche essere negativa}, ma in quel caso indica che cambia il verso.
	\subsection{Seconda legge di Ohm}
	\begin{equation}
		R = \rho \cdot \frac{l}{S} \;\; [\Omega \cdot m]
	\end{equation}
	$R$ indica la resistenza \\
	$\rho$ indica la \textbf{resistività}, ovvero quanto un materiale resiste all'elettricità, $\rho \in \left[ 10^{-7}, 10^{16}\right]$ \\
	$l$ indica la \textbf{lunghezza} del campione in cui scorre la corrente \\
	$S$indica la \textbf{sezione/area} del campione in cui scorre la corrente \\ \\
	L'inverso della resistività è la \textbf{conducibilità}, ovvero quanto un materiale conduce elettricità, viene indicata con:
	\[ \sigma = \frac{1}{\rho} \]
	\subsubsection{Variazione di Resistività (o Concudibilità) con la temperatura}
	\subsection{Bipoli}
	\subsubsection{Bipolo Elettrico (generico)}
	\subsubsection{Filo Elettrico}
	\subsubsection{Bipolo "Resistenza"}
	\subsubsection{Bipolo "Generatore Ideale di Tensione"}
	\subsubsection{Bipolo "Generatore Ideale di Corrente"}
	\subsection{Convenzione del Generatore}
	Convenzione in cui la tensione $\Delta V$ e la corrente $I$ elettrica hanno lo stesso verso, in questo caso la prima legge di Ohm diventa:
	\begin{equation}
		\Delta V = - (R \cdot I)
	\end{equation}
	\subsection{Convenzione del'Utilizzatore}
	Convenzione in cui la tensione $\Delta V$ e la corrente $I$ elettrica hanno verso discorde, in questo caso la prima legge di Ohm diventa:
	\begin{equation}
		\Delta V = R \cdot I
	\end{equation}
	\subsection{Principio di Kirchhoff al Nodo (1° Principio di Kirchhoff)}
	\subsection{Principio di Kirchhoff alla Maglia (2° Principio di Kirchhoff)}
	\subsection{Conduttanza Elettrica}
	\subsection{Potenza Elettrica di un Bipolo}
	\subsubsection{Convenzione Utilizzatore}
	\subsubsection{Convenzione Generatore}
	\subsection{Fenomeni Termici}
	\subsubsection{Legge di Ohm Termica}
	\subsubsection{Resistenza Termica}
	\subsection{Collegamenti fra Bipoli}
	\subsubsection{Collegamento in Serie}
	\subsubsection{Collegamento in Parallelo}
	\subsubsection{Partitore di Tensione (Serie di Resistenze)}
	\subsubsection{Partitore di Corrente (Parallelo di Resistenze)}
	\subsubsection{Serie di Generatori Ideali di Tensione}
	\subsubsection{Parallelo di Generatori Ideali di Tensione}
	\subsubsection{Serie di Generatori Ideali di Corrente}
	\subsubsection{Parallelo di Generatori Ideali di Corrente}
	\subsection{Generatori Reali}
	\subsubsection{Generatore Reale di Tensione}
	\subsubsection{Generatore Reale di Corrente}
	\subsubsection{Equivalenze fra Generatori Reali}
	\subsection{Generatori Reali}
	\subsubsection{Generatore Reale di Tensione}
	\subsubsection{Generatore Reale di Corrente}
	\subsubsection{Equivalenze fra Generatori Reali}
	\subsubsection{Considerazioni Energetiche sui Generatori Reali}
	\subsubsection{Generatori Reali di Tensione in Serie}
	\subsubsection{Generatori Reali di Corrente in Serie}
	\subsection{Strumenti di Misura}
	\subsubsection{Voltmetro}
	\subsubsection{Amperometro}
	\subsubsection{Ohmero}
	\subsubsection{Wattmetro}
	\subsection{Generatori Pilotati}
	\subsubsection{Generatore di Tensione Pilotato da Una Tensione}
	\subsubsection{Generatore di Tensione Pilotato da Una Corrente}
	\subsubsection{Generatore di Corrente Pilotato da Una Tensione}
	\subsubsection{Generatore di Corrente Pilotato da Una Corrente}
	\subsection{Teorema di Millman}
	\subsection{Trasformazioni Stella-Triangolo}
	\subsection{Principio di Sovrapposizione degli Effetti}
	\subsubsection{Circuiti Lineari}
	\subsubsection{Metodo della Sovrapposizione degli Effetti}
	\subsection{Teorema di Thevenin}
	\subsection{Teorema di Norton}
	\subsection{Thevenin e Norton in Presenza di Generatori Pilotati}
	\subsection{Metodi Risolutivi}
	\subsubsection{Metodo delle Correnti di Maglia}
	\subsubsection{Metodo dei Potenziali di Nodo}
	\subsubsection{Casi Degeneri}
\end{document}